\documentclass{article}
\title{An Automap for Essays}
\begin{document}
\begin{abstract}
What the Descent automap does well, where its late levels turn into a thicket, and the rules an essay automap should follow instead.
\end{abstract}
\section{Where Descent's Map Succeeds}\label{sec:map-success}
In the simpler levels the wireframe gives an immediate topology: a central chamber, branches, a few loops, recognisable ends. In the dense late levels, drawing every edge at once becomes a thicket. The map stays usable there only because the pilot has already learned the space by moving through it.
\section{Egocentric Before Exhaustive}\label{sec:map-ego}
% mine: thesis
The map should open in the reader's current orientation, centred on the present chamber, with the route just travelled highlighted and remote geometry dimmed. A local map with seven recognisable connections is easier to use than a truthful picture of a whole corpus.
\section{Sparse, Invariant Colour}\label{sec:map-colour}
White carries the skeleton. Colour marks what a reader can act on: the current route, exits to other works, prerequisite doors, open threads, beacons and the present chamber. Illustrations appear on the map only as simple landmark glyphs.
\section{Three States of Discovery}\label{sec:map-states}
Unexplored structure is absent. Revealed but unvisited structure is a faint envelope. Traversed passages are solid. A brief scan may show more for a moment, but only being there makes it permanent.
% thread[open]: whether revealed structure should ever fade again is left open
\section{Semantic Depth}\label{sec:map-depth}
% mine: exit=two-games#sec:games-loops
Viewing distance should change the representation, not only its size. Close in, the map shows the current section and its alcoves; further out, the part; further still, the whole work; beyond that, the work becomes one system among others. Geometry belongs only at the scale where it affects the next decision.
\end{document}
